\documentclass[a4paper,12pt]{article}
\setlength{\headheight}{68.803pt}
\addtolength{\topmargin}{-56.803pt}

\usepackage[utf8]{inputenc}
\usepackage[T1]{fontenc}
\usepackage{lmodern}
\usepackage[nopatch=footnote]{microtype}
\usepackage{fvextra}
\DefineVerbatimEnvironment{Verbatim}{Verbatim}{
  breaklines=true,
  breakanywhere=true,
  frame=single,
  fontsize=\small
}

\usepackage[
  left=1in,
  right=1in,
  top=3cm,
  bottom=1in
]{geometry}

\usepackage[
    colorlinks=true,
    linkcolor=black,
    urlcolor=blue,
    citecolor=blue,
    pdfborder={0 0 0},
    pdfauthor={Massimo Mantineo},
    pdftitle={HO17: Distributed Bellman-Ford},
    pdfsubject={Laboratorio di Reti e Sistemi Distribuiti: HANDS-ON 17},
]{hyperref}

\usepackage{graphicx}
\graphicspath{{../../assets/}}
\usepackage{tikz}
\usetikzlibrary{positioning, arrows.meta, shapes.symbols}
\usepackage{listings}
\usepackage{xcolor}
\usepackage{amsmath, amssymb}
\usepackage{fancyhdr}
\usepackage{setspace}
\usepackage{pgfplots}
\pgfplotsset{compat=1.17}
\usepackage{booktabs}
\usepackage[skins, listings]{tcolorbox}
\usepackage{float}

\newtcblisting{terminal}[1][]{
    listing only,
    colback=black!85!white,
    colframe=black!70!white,
    colupper=green!80!white,
    coltitle=white,
    title=#1,
    fonttitle=\bfseries\sffamily,
    listing options={
        language={},
        basicstyle=\ttfamily\small\color{green!80!white},
        backgroundcolor={},
        frame=none,
        breaklines=true,
        showstringspaces=false,
        numbers=none
    },
    enhanced,
    drop shadow,
    arc=3mm
}

\newtcblisting{ccode}[1][]{
    listing only,
    colback=blue!3!white,
    colframe=blue!70!black,
    title=#1,
    fonttitle=\bfseries\sffamily,
    listing options={
        style=cstyle,
        backgroundcolor={},
        frame=none
    },
    enhanced,
    drop shadow,
    arc=3mm
}

\setlength{\footskip}{50pt}

\lstdefinestyle{cstyle}{
    language=C,
    basicstyle=\ttfamily\small,
    numbers=left,
    numberstyle=\tiny,
    stepnumber=1,
    numbersep=5pt,
    backgroundcolor=\color{white},
    showspaces=false,
    showstringspaces=false,
    showtabs=false,
    frame=single,
    tabsize=4,
    captionpos=b,
    breaklines=true,
    breakatwhitespace=true,
    keywordstyle=\color{blue}\bfseries,
    commentstyle=\color{green!50!black},
    stringstyle=\color{red},
}
\lstset{style=cstyle}

\pagestyle{fancy}
\fancyhf{}

\fancyhead[L]{%
  \parbox[b]{0.45\textwidth}{%
    \small
    Student Name: Massimo Mantineo\\
    Student ID: 541924
    \vspace{20pt}
  }%
}
\fancyhead[C]{%
  \parbox[b]{0.2\textwidth}{%
    \centering
    \normalsize \bfseries
    LabReSiD25\\
    Hands-On 17
    \vspace{20pt}
  }%
}
\fancyhead[R]{%
  \parbox[b]{0.35\textwidth}{%
    \flushright
    \includegraphics[width=3.5cm]{\detokenize{logo_unime_orizontale.png}}
    \vspace{15pt}
  }%
}

\renewcommand{\contentsname}{Indice}

\title{\vspace{-2cm}Report (HO17)\\
\large Algoritmo Bellman-Ford: Distributed Shortest Path via IPC}
\author{Massimo Mantineo \\ Matricola: 541924}
\date{MESSINA, A.A. 2024/2025}

\begin{document}

\begin{titlepage}
    \thispagestyle{empty}
    \centering
    \vspace*{1cm}
    \includegraphics[width=6cm]{\detokenize{logo_unime_orizontale.png}}\par\vspace{1cm}
    \vspace{2cm}
    {\Large \textbf{Università degli Studi di Messina}\par}
    {\large Dipartimento MIFT - Corso di Laurea Triennale in Informatica (L-31)\par}
    \vspace*{1cm}
    {\large Laboratorio di Reti e Sistemi Distribuiti\par}
    \vspace{1cm}
    {\Large \textbf{HANDS-ON 17:}\\
    Algoritmo di Bellman-Ford per il Distance-Vector Routing\par}
    \vspace{2cm}
    {\large Massimo Mantineo \par}
    {\large Matricola: 541924 \par}
    \vfill
    {\large Messina, A.A. 2024/2025 \par}
\end{titlepage}

\clearpage
\pagestyle{fancy}
\tableofcontents
\newpage

\section{Problem Statement}
Nell'ambito dell'astrazione logica di nodi di rete come processi UNIX, si richiede l'implementazione del celebre algoritmo matematico di \textbf{Bellman-Ford} per il monitoraggio dei grafi pesati. Si tratta dell'archetipo fondante per la classe di protocolli \textit{Distance-Vector} usati dai router fisici per eleggere la minima distanza hop-by-hop verso una sottorete (Si veda, ad esempio, lo standard `RIP` o `EIGRP`). Ancora una volta l'elaborazione non sarà accentrata su un database monolitico for-loop classico, ma si propagherà in maniera genuinamente \textbf{Distribuita ed Asincrona} sfruttando \texttt{fork} ed appoggiandosi alle Syscall \texttt{mkfifo} per le code di messaggistica.

\section{Stato dell'Arte}
L'algoritmo di Bellman-Ford assicura, in reti prive di cicli a costo negativo, la convergenza totale dei pesi. 
Esso si esprime tramite l'assioma di calcolo iterativo:
\begin{equation}
    D(x, y) = \min_v \left\{ c(x,v) + D(v,y) \right\}
\end{equation}
Dove:
\begin{itemize}
    \item $D(x,y)$: Il valore stimato attualizzato della distanza dal nodo $x$ al target $y$.
    \item $c(x,v)$: Il peso prefissato nel cavo (Edge) che separa il nodo adiacente $x$ dal nodo limitrofo $v$.
\end{itemize}

Ogni nodo \textit{riceve} in asincrono messaggi in buffer dalla propria FIFO, decodifica il Vettore $D(v,y)$ ricevuto, somma all'overhead biologico $c(x,v)$ che esso prova per raggiungere $v$, e decreta se il cammino calcolato è più economico di quello che custodiva in memoria, mutando ed allertando i vicini di rimando in caso di Update positivo.

\section{Metodologia ed Implementazione}

Siccome il grafo visivo manca all'appello, optiamo per disporre una trincea a \textbf{Cost Matrix Asimmetrica} per 4 Nodi:
\begin{ccode}
// Single-Source cost_matrix
// Node 0 is the Target.
int cost_matrix[NUM_NODES][NUM_NODES] = {
    {0, 2, 5, INF}, // Edge out dal nodo 0
    {2, 0, 1, 4},   // Edge out dal nodo 1
    {5, 1, 0, 1},   // Edge out dal nodo 2
    {INF, 4, 1, 0}  // Edge out dal nodo 3
};
\end{ccode}
Abbiamo programmato volontariamente un trabocchetto: benché esista l'arco diretto Nodo 2 $\rightarrow$ Target 0 al costo di asticella `5`, la somma dei pesi passando attivamente per Nodo 1 costa solamente `3` ($1+2$). Costringiamo così in maniera fittizia l'algoritmo a ricalcolare le metriche ignorando l'arco apparentemente più palese (Il nodo 3 dista `1` dal 2, per un totale di `4`, ed in quanto esso dista `4` dal Nodo 1 che aggiunge il fatidico `2`, la via più ovvia peserebbe `6`. Viene scelto `4`).


\subsection{Istruzioni di Esecuzione}
L'applicativo viene gestito interamente tramite Makefile:
\begin{terminal}[Build & Run]
$ make
$ ./programma
\end{terminal}

\section{Risultati}
Si è provveduto allo sfoggio del binario, compilato per iterazioni continue ed abortito dopo $3_{sec}$ dall'incessazione delle scoperte IPC (Timeout Convergenza).

\begin{terminal}[Mutua Esclusione Distribuita - Output]
massimo@PORTATILE-Linux:~/Scrivania/LabReSid24-25/hands-on/ho17$ ./bellman_ford 
===== Bellman-Ford Algoritmo di Distance Vector Distribuito =====
Obiettivo: Cammini ottimali da ogni Nodo verso il NODO TARGET (0)

[NODO 0 - TARGET] Distanza zero locale. Segnalo Inizio alle adiacenze.
[NODO 1] D-V UPDATE! Vett. 999 > 2 (Nuovo hop: Nodo 0). Espando l'onda d'urto...
[NODO 2] D-V UPDATE! Vett. 999 > 5 (Nuovo hop: Nodo 0). Espando l'onda d'urto...
[NODO 3] D-V UPDATE! Vett. 999 > 6 (Nuovo hop: Nodo 1). Espando l'onda d'urto...
[NODO 2] D-V UPDATE! Vett. 5 > 3 (Nuovo hop: Nodo 1). Espando l'onda d'urto...
[NODO 3] D-V UPDATE! Vett. 6 > 4 (Nuovo hop: Nodo 2). Espando l'onda d'urto...

--- [NODO 2 ESTINTO/CONVERSO] Shortest Path Totale => 3 (via Nodo 1) ---
--- [NODO 0 ESTINTO/CONVERSO] Shortest Path Totale => 0 (via Nodo 0) ---
--- [NODO 3 ESTINTO/CONVERSO] Shortest Path Totale => 4 (via Nodo 2) ---
--- [NODO 1 ESTINTO/CONVERSO] Shortest Path Totale => 2 (via Nodo 0) ---

==== Convergenza Assoluta Raggiunta. Discesa IPC ====
\end{terminal}

\section{Conclusioni e Sviluppi Futuri}
La traccia IPC e i log confermano clamorosamente le nostre predisposizioni da scrivania. All'insorgere del tempo e con l'avanzamento dei cicli, il **NODO 2 e il NODO 3 scartano le rotte precedenti (5 e 6 rispettivamente)** per abbracciare l'Update di vettori asincrono che segnalava la presenza del percorso a costo calcolato `3` e `4`, dimostrando ad amplifero la correttezza del demone C. Le sysfile temporanee di buffer (Pipe multiple) vengono annibillate dolcemente dallo sterrato in chiusura dell'ambiente Unix locale.

\end{document}
